% !TEX TS-program = pdflatex
\documentclass[conference]{IEEEtran}
\usepackage{cite}
\usepackage{amsmath,amssymb,amsfonts}
\usepackage{mathtools}
\usepackage{algorithmic}
\usepackage{graphicx}
\usepackage{textcomp}
\usepackage{booktabs}
\usepackage{multirow}
\usepackage[table]{xcolor}

\def\BibTeX{{\rm B\kern-.05em{\sc i\kern-.025em b}\kern-.08em
    T\kern-.1667em\lower.7ex\hbox{E}\kern-.125emX}}

\begin{document}

\title{Self Distillation using only supervised fine tuning: a duplex cascade case study}

\author{
\IEEEauthorblockN{Penh Fel}
\IEEEauthorblockA{}
}

\maketitle

\begin{abstract}
Fine-tuning a full-duplex speech model to emit conversational turn-taking tags
(the \emph{duplex cascade} micro-turn protocol) with plain supervised fine-tuning
(SFT) regresses the model's general capabilities: function calling, instruction
following and world knowledge degrade because the off-policy training signal
pushes the whole next-token distribution toward tag-heavy dialogue. We show that
this regression can be avoided \emph{without reinforcement learning and without
on-policy rollouts} by borrowing the self-distillation idea of
\emph{Self-Distillation Fine-Tuning} (SDFT): the frozen original model acts as its
own teacher, and the fine-tune optimizes a two-term loss---the weighted
cross-entropy on the duplex tags exactly as before, plus a per-token
\textbf{forward-KL to the frozen base model} on every other (content) position.
The content term \emph{anchors} the student to the pretrained distribution
everywhere the tags are not explicitly supervised, which is precisely the
mechanism that prevents regression. We train a QLoRA adapter on 10k teacher-derived
duplex dialogues and evaluate against both the base model and the plain-SFT
baseline. The distill model learns the protocol (96.4\% tag accuracy) while keeping
function calling and general instruction-following at base level (100\% vs the SFT
baseline's 55.6\% and 41.7\%), and it stays within 1--2 pp of base on MMLU, HellaSwag,
ARC and TruthfulQA, and above base on IFEval---where the SFT baseline drops broadly.
The only measured regression is on multi-step math reasoning (GSM8K). These results
establish per-token KL-anchoring to the frozen base as a practical, RL-free
anti-forgetting mechanism for task-specific fine-tuning.
\end{abstract}

\begin{IEEEkeywords}
self-distillation, continual learning, supervised fine-tuning, full-duplex dialogue,
turn-taking, catastrophic forgetting, function calling
\end{IEEEkeywords}

\section{Introduction}
\IEEEPARstart{L}{arge} language models that power interactive systems are often
post-trained on narrowly defined behavioral protocols. A representative example is
the \emph{duplex cascade} paradigm~\cite{duplexcascade}: a conversational LLM is
fine-tuned to emit special turn-taking tokens (``\textless|user is speaking|>'',
``\textless|user finish speaking|>'', ``\textless|user interruption|>'', etc.) that
drive a VAD-free, full-duplex speech pipeline. Learning the protocol is easy---the
supervision is cheap, synthetic and deterministic---but the fine-tune is destructive.
In our pt-BR implementation of the duplex cascade, the plain-SFT model that learns
the tags with 98\% accuracy simultaneously loses most of its ability to call
functions (100\% $\rightarrow$ 55.6\%) and follow general instructions
(100\% $\rightarrow$ 41.7\%). This is catastrophic forgetting in its most practical
form: the model becomes a good turn-taker and a bad assistant.

The conventional remedy is to make the post-training signal \emph{on-policy}, via
reinforcement learning (RL) with an explicit reward~\cite{chu2025,airef}. But rewards
are often unavailable, and on-policy RL (PPO/GRPO) is expensive and brittle. A recent
line of work, \emph{Self-Distillation Fine-Tuning}
(SDFT)~\cite{selfdistill2026}, obtains the benefits of on-policy learning from
demonstrations alone: the model is used in two roles---a \emph{teacher} conditioned
on a demonstration, and a \emph{student} conditioned only on the query---and the
training minimizes a KL divergence between the two on the student's own trajectories.
SDFT substantially reduces catastrophic forgetting compared to SFT, but it still
requires on-policy rollouts (roughly $4\times$ the wall-clock of SFT).

We show that for the duplex-cascade setting there is a much simpler instantiation
that needs \emph{neither RL nor rollouts}. Our key observation is that the
``new skill'' (the turn-taking protocol) and the ``everything else'' (the general
capabilities we want to preserve) are \emph{separable by token position}: the tags
live at supervised positions, the content lives at the other positions. We therefore
split the objective in two:

\begin{itemize}
  \item \textbf{Tags}: the weighted cross-entropy to the supervised tag tokens,
        exactly as in the plain SFT recipe.
  \item \textbf{Everything else}: a per-token forward-KL to the \emph{frozen base
        model} on content positions, so the student's distribution is anchored to
        the pretrained policy wherever there is no explicit tag supervision.
\end{itemize}

The teacher is the frozen original model itself (conditioned on the same duplex
sequence via the tag tokens in its input), so no second model and no rollouts are
needed. Because the student is \emph{base + LoRA}, the teacher forward is obtained
with zero extra VRAM by temporarily disabling the adapters. We precompute the
teacher's top-$k$ next-token logits once (offline cache), which removes the second
forward pass entirely and makes the run no more expensive than SFT.

We evaluate on a pt-BR duplex-cascade model (Qwen3-4B-Instruct-2507, QLoRA, 10k
teacher-derived dialogues). The resulting \emph{distill} model matches the plain-SFT
model on the protocol (96.4\% vs 98\% tag accuracy) while retaining base-level
function calling, general instruction following, MMLU, HellaSwag, ARC and TruthfulQA,
and \emph{improving} on IFEval. The plain-SFT baseline regresses on all of them. The
one caveat is GSM8K, where the distill model drops 11 pp below base---a limitation we
discuss in Section~\ref{sec:discussion}.

\section{Background and Related Work}
\subsection{Catastrophic forgetting and off-policy fine-tuning}
Sequential SFT is inherently off-policy: the model learns to imitate expert actions
on a fixed offline distribution, and prior work has shown this leads to severe
forgetting when the model is adapted to a new task or domain~\cite{kirkpatrick2017,
lifhoiem2017,mecklenburg2024}. On-policy methods---training on trajectories induced
by the model itself---forget substantially less~\cite{ross2011,airef,chu2025,
rlsrazor2025}. DAgger~\cite{ross2011} established that on-policy imitation avoids the
compounding-error failure mode of behavior cloning. More recently,
``SFT memorizes, RL generalizes''~\cite{chu2025} and related analyses~\cite{airef}
showed empirically that RL-based post-training retains general capabilities far
better than SFT.

\subsection{Self-distillation}
Our method is most closely related to \emph{Self-Distillation Fine-Tuning}
(SDFT)~\cite{selfdistill2026}. SDFT builds the teacher by conditioning the model on
an expert demonstration, $\pi(\cdot|x,c)$, and trains the student $\pi_\theta(\cdot|x)$
to match it on the student's own rollouts:
\begin{equation}
\begin{split}
  L(\theta) \;=\; D_{KL}\big(\pi_\theta(\cdot|x) \,\|\, \pi(\cdot|x,c)\big)\\
  \;=\; \mathbb{E}_{y\sim\pi_\theta}\sum_{t}
      \log\frac{\pi_\theta(y_t|y_{<t},x)}
      {\pi(y_t|y_{<t},x,c)}\,.
\end{split}
\end{equation}
The paper shows that on Tool Use, Science Q\&A and Medical skill-learning, SDFT
achieves higher new-task accuracy than SFT while keeping previous-capabilities
performance near the base model, and that an EMA of the student is the most stable
teacher. Our work differs in two ways. First, we do \emph{not} run on-policy
rollouts; instead the tags are supplied by a deterministic construction rule, so the
only on-policy-sensitive component is the content, which we anchor to the frozen base
distribution directly. Second, our teacher is the frozen base model (not a
demonstration-conditioned variant), which is exactly the ``keep everything else the
same'' reference the anti-forgetting objective requires.

Context distillation~\cite{bai2022,snell2022} also uses a model conditioned on extra
context as a teacher for the model without it, but typically distills offline from
static contexts and does not combine a separate supervised tag term with a
content-anchoring KL. \emph{Learning by distilling context}~\cite{snell2022}
distills a prompt-conditioned teacher into the base model; our tag term is the
analogous ``hard'' supervision and our content term the distillation.

\subsection{Parameter-efficient fine-tuning}
We train with QLoRA~\cite{qlora2023} (4-bit NF4 base, LoRA~\cite{lora2022} rank 16,
$16$ $\alpha$, dropout 0) via the Unsloth~\cite{unsloth} training library, keeping the
embedding and LM head trainable so the new special-token rows can be learned.

\section{Method: DuplexCascade-Distill}
\label{sec:method}
\subsection{Setup}
Let the base model be $\pi_0$ (frozen) and the student $\pi_\theta$ be the same
weights plus a LoRA adapter ($\theta$ = adapter parameters). The training corpus is a
set of \emph{duplex sequences} $x$ produced by the duplex micro-turn builder: a
dialogue is split into micro-turns and conversational special tokens
($\mathcal{T}$) are inserted deterministically at turn boundaries
(e.g. ``\textless|user finish speaking|>'' before every system turn, probabilistic
interruptions/backchannels). Every position $t$ of $x$ is thus either a
\emph{tag} position (label $y_t \in \mathcal{T}$) or a \emph{content} position
(label $y_t$ is ordinary text). User-side and header positions are masked
(label $-\infty$), as in the original recipe.

\subsection{Teacher content}
The tags themselves come from the deterministic builder. The \emph{content} comes
from the teacher: the frozen base model, prompted with a \emph{duplex-cascade trigger}
system message (``you are a DuplexCascade full-duplex voice assistant\ldots speak in
short, natural micro-turns\ldots no tags''), generates each assistant turn
conditioned on the history. At first the base model does not follow the protocol
faithfully---that is expected and irrelevant; we keep only its \emph{content} and
let the deterministic builder place the tags. This yields the aligned demonstration
sequence $x$.

\subsection{Loss}
Let $T$ be the set of tag positions and $C$ the set of content positions. The
student loss is
\begin{equation}
\label{eq:loss}
\begin{split}
  L(\theta) \;=\;
  \underbrace{\sum_{\mathclap{t\in T}} w_t\, \mathrm{CE}\big(\pi_\theta(\cdot|x_{<t}),\, y_t\big)
  }_{\text{tags: as in plain SFT}}\\
  \qquad{}+{}\;
  \lambda\underbrace{\sum_{\mathclap{t\in C}} D_{KL}\big(
      \pi_0(\cdot|x_{<t}) \,\|\, \pi_\theta(\cdot|x_{<t})\big)
  }_{\text{content: anchor to the frozen base}},
\end{split}
\end{equation}
where $w_t$ are the per-special-token weights of the original recipe (e.g.
``\textless|user finish speaking|>''$=10$, ``\textless|user interruption|>''$=5$).
The KL is in the forward direction and computed in cross-entropy form
$-\sum_k p_0(k)\log p_\theta(k)$ over the teacher's top-$k$ ($k{=}32$) next-token
support, with the student logits gathered at those ids. Because the teacher is
conditioned on the same $x$ (its input already contains the tags), its content
distribution is the natural ``what would the original model say in a duplex
dialogue'' reference. The single hyper-parameter $\lambda$ trades tag fidelity
against anchoring; we use $\lambda{=}1.0$.

\subsection{Practical implementation}
\begin{itemize}
\item \textbf{Teacher forward is free.} The student is base $+$ LoRA, so the teacher
      is the same weights without adapters: one additional \texttt{no\_grad} forward
      under \texttt{disable\_adapter()} costs zero extra VRAM.
\item \textbf{Offline teacher cache.} To avoid the second forward in the training
      loop (and the GPU-allocator pressure it creates at the 16 GB ceiling), we
      precompute the teacher's top-$k$ ids/logits for every position of every
      sequence once (one bf16 forward pass over the corpus, $\approx$36 min). Training
      then reads cached targets and runs at the plain-SFT cost. This is the
      difference between our ${\sim}10$ s/step and the two-forward path.
\item \textbf{Anchor set (D2 variant).} A small set of general-capability prompts
      (function calling, QA) is distilled the same way, pinning the LoRA delta to
      $\approx 0$ outside the duplex protocol.
\item \textbf{Artifact masking.} Following the SDFT ``learned artifacts''
      fix~\cite{selfdistill2026}, the first few content tokens after each
      ``\textless|user finish speaking|>'' tag are left unsupervised so the student
      does not copy teacher surface artifacts.
\end{itemize}

\section{Experimental Setup}
\label{sec:setup}
\subsection{Models}
Base: \texttt{Qwen/Qwen3-4B-Instruct-2507} (dense Qwen3, 151.9k vocab). Student:
QLoRA (r$=16$, $\alpha{=}16$, dropout 0, all linear modules, trainable embeddings and
LM head), 7 duplex special tokens appended and Gaussian-initialized. Teacher: the same
base, frozen. All models evaluated as bf16.

\subsection{Data}
10,000 pt-BR dialogues from the existing corpus had every assistant turn regenerated
by the teacher (trigger prompt, temperature 0.7, top-p 0.95), then passed through the
deterministic micro-turn builder (user chunks 1--7 tokens, system chunks 10--48,
probabilistic pauses/interruptions/backchannels), yielding 10,249 training sequences
at ${\le}1024$ tokens. 300 anchor prompts (150 function calling, 150 general QA).

\subsection{Training}
2,000 optimizer steps, batch 1 $\times$ grad-accum 16 (effective 16), AdamW lr
$10^{-4}$ (embeddings $2{\times}10^{-4}$), 20-step warmup, bf16, gradient
checkpointing, $\lambda{=}1.0$ (chosen by a short sweep on $\lambda\in\{0.1,0.3,1.0\}$
using tag accuracy and KL drift). The \texttt{sft} baseline is the existing
DuplexCascade-PT QLoRA fine-tune (same data family, no KL term).

\subsection{Evaluation}
\emph{Protocol}: tag accuracy (argmax $=$ supervised tag) on a held-out duplex eval
split. \emph{Capabilities}: 18 pt-BR function-calling prompts (native Qwen tool
calls, greedy) and 12 general instruction/QA prompts. \emph{Distribution drift}:
mean per-token forward-KL (student $\|$ frozen base) on the content positions of the
held-out split---a direct measure of ``everything else stayed the same''.
\emph{Standard benchmarks}: HellaSwag, ARC-Challenge, MMLU (2000-example subsets),
GSM8K (5-shot, 300), TruthfulQA (mc2), IFEval (40) via lm-evaluation-harness, raw
prompts, no chat template, identical configs for all three models.

\section{Results}
\label{sec:results}
\subsection{Protocol vs. capability retention}
Table~\ref{tab:custom} shows the central result. The \emph{distill} model learns the
protocol (96.4\% tag accuracy, only 1.6 pp behind the SFT baseline) while keeping
function calling and general instruction-following at \emph{base level} (100\%), and
a KL drift only slightly above the base floor. The plain-SFT baseline, by contrast,
learns the protocol (98\%) but regresses catastrophically on the format-following
capabilities (function calling 55.6\%, general 41.7\%) and drifts ${\sim}3.5\times$
further from the base distribution.

\begin{table}[t]
\centering
\caption{Protocol and capability retention (our pt-BR eval sets).}
\label{tab:custom}
\begin{tabular}{lccc}
\toprule
 & base & \textbf{distill} & sft \\
\midrule
Tag accuracy (\%) & 0.0 & \textbf{96.4} & 98.0 \\
Function calling (\%) & 100 & \textbf{100} & 55.6 \\
General instruction/QA (\%) & 100 & \textbf{100} & 41.7 \\
KL drift vs base (nats) & 0.51 & 1.07 & 4.76 \\
\bottomrule
\end{tabular}
\end{table}

\subsection{Standard benchmarks}
Table~\ref{tab:bench} reports the lm-evaluation-harness results. The distill model is
within ${\sim}1$--2 pp of base on knowledge (MMLU), commonsense (HellaSwag) and
reasoning (ARC), \emph{above} base on factuality (TruthfulQA $+2.3$ pp) and
instruction following (IFEval loose $+3.2$ pp), and at base on the format-following
evals of Table~\ref{tab:custom}. The SFT baseline drops on MMLU ($-4.5$ pp),
TruthfulQA ($-2.5$ pp), IFEval ($-12.7$ pp) and the custom sets. The one regression
of the distill model is GSM8K ($-11.3$ pp), discussed next.

\begin{table}[t]
\centering
\caption{Standard benchmarks (lm-evaluation-harness 0.4.12; subsets as noted).}
\label{tab:bench}
\footnotesize
\resizebox{\columnwidth}{!}{%
\begin{tabular}{lccc}
\toprule
Benchmark & base & \textbf{distill} & sft \\
\midrule
HellaSwag acc-norm ($n{=}2000$) & 0.584 & 0.566 & 0.593 \\
ARC-Challenge acc-norm ($n{=}2000$) & 0.584 & 0.563 & 0.574 \\
MMLU ($n{=}2000$) & 0.705 & 0.698 & \textbf{0.660} \\
GSM8K exact ($n{=}300$, 5-shot) & 0.890 & 0.777 & 0.827 \\
TruthfulQA mc2 & 0.626 & \textbf{0.649} & 0.601 \\
IFEval loose / strict ($n{=}40$) & 0.730/0.651 & \textbf{0.762/0.683} & 0.603/0.524 \\
\bottomrule
\end{tabular}%
}
\end{table}

\section{Discussion and Limitations}
\label{sec:discussion}
\emph{Why the content-KL works.} The tags live at supervised positions and are
predicted with the original high weights, so they are learned essentially as before.
Every other position is anchored to the frozen base, which is exactly the behavior
we want to preserve. The mechanism is a single-stage, per-token analogue of the
``re-invoke'' two-stage procedure~\cite{tml2025} and of SDFT's on-policy
anchoring~\cite{selfdistill2026}, but it requires neither rollouts nor a reward.

\emph{The GSM8K caveat.} The distill model drops 11 pp below base on GSM8K, a
long-form multi-step reasoning task. This is the honest tradeoff of a strong
anchoring term ($\lambda{=}1.0$): content distributions are heavily pinned to the
base, and on 300 samples the 4B base's fragile chain-of-thought is sensitive to any
distributional shift. Two natural follow-ups are (i) a lower $\lambda$ (our sweep did
not include a math benchmark), and (ii) the EMA-teacher or on-policy variants of SDFT
(SPEC D3/D4), which track the student's progress and may preserve the student's
adaptability on reasoning while still reducing forgetting.

\subsection{Next steps: learning only part of the answer}
The crux of our method is the \emph{positional separation} of the objective: a small,
well-localized set of supervised positions carries the new behavior, and every other
position is anchored to the frozen base. This is a special case of a much more common
and important fine-tuning pattern---\emph{partial-answer supervision}---where the
student's target output is already largely correct and only a small, identifiable
part must be edited. Our approach should transfer directly to any such setting, and
we plan to evaluate it there next. Concretely:

\begin{itemize}
  \item \textbf{Guardrail classifiers.} The most immediate target is fine-tuning an
        LLM to behave as a \emph{classifier}. For a content-moderation guardrail, the
        input is a text and the target output is a verdict from a small fixed
        taxonomy (e.g.\ ``\textsc{ok}'' / ``\textsc{flagged}'', or a category label),
        often preceded by a one-line justification. Plain SFT on this data collapses
        the model's output distribution onto the label tokens and erodes general
        abilities. With our recipe, the \emph{verdict position is the ``tag''} (weighted
        cross-entropy) and the justification and everything else are the ``content''
        (forward-KL to the frozen base). The model then learns the decision boundary
        while keeping its language, instruction-following and knowledge intact---exactly
        the retention we measured for the duplex tags (Table~\ref{tab:custom}).
  \item \textbf{Other candidate settings.} (i)~\emph{Schema/format fix-ups} where the
        answer semantics are correct but the formatting (JSON keys, tool-call syntax)
        must change; (ii)~\emph{claim correction} (medical, legal, RAG) where only a
        factually wrong token span needs editing; (iii)~\emph{code repair} where only
        the buggy lines change; (iv)~\emph{citation insertion}.
\end{itemize}

The enabling requirement is that the ``new behavior'' is locatable at predictable
token positions. When it is---a verdict token, a tag, a format delimiter---the
off-policy recipe above is directly applicable and stays at plain-SFT cost. When the
edit location is \emph{not} known a priori, one can first align the student to a
canonical template (exactly as the deterministic micro-turn builder does for duplex
dialogue) so the edit collapses onto predictable positions; otherwise the on-policy
SDFT variant~\cite{selfdistill2026} is the fallback. This planned evaluation also
speaks to the principle of \emph{RL's Razor}~\cite{rlsrazor2025}: forgetting is
governed by the KL to the base on the new task, and partial-answer settings are
precisely where an explicit per-token KL constraint is cheapest and most natural,
because the supervised support is small.

\emph{Limitations.} Our evaluation is on a single backbone (Qwen3-4B), a single
language (pt-BR) and subset-scale benchmarks; the IFEval and GSM8K numbers carry
per-sample stderr of order $2$ pp. We did not ablate $\lambda$ against math, nor test
the full 50k-scale data regime. The offline cache assumes the teacher is fixed,
which is true by construction but means the student can never ``teach forward'' a
better protocol---acceptable here because the tags are deterministic.

\section{Conclusion}
We showed that the duplex-cascade protocol can be learned with plain supervised
fine-tuning \emph{without} the catastrophic regression that usually accompanies it,
by splitting the objective into a supervised tag term and a per-token forward-KL to
the frozen base model. The approach needs neither reinforcement learning nor on-policy
rollouts, and with an offline teacher cache it costs no more than SFT. On a pt-BR
duplex-cascade model it matches the plain-SFT protocol accuracy while preserving
base-level function calling, general instruction following, MMLU, HellaSwag, ARC and
TruthfulQA and improving IFEval---where the SFT baseline regresses broadly. Per-token
KL-anchoring to the frozen base is a practical, RL-free anti-forgetting tool that
generalizes beyond duplex dialogue.

\begin{thebibliography}{99}

\bibitem{selfdistill2026}
I.~Shenfeld, M.~Damani, J.~H{\"u}b{\"o}tter, and P.~Agrawal,
``On-policy self-distillation fine-tuning (SDFT): self-distillation enables continual
learning,'' arXiv:2601.19897, 2026.

\bibitem{duplexcascade}
J.~Yang, Y.~Fujita, and Y.~Sudo,
``DuplexCascade: Full-duplex speech-to-speech dialogue with VAD-free cascaded
ASR--LLM--TTS pipeline and micro-turn optimization,'' arXiv:2603.09180, 2026.

\bibitem{ross2011}
S.~Ross, G.~Gordon, and D.~Bagnell,
``A reduction of imitation learning and structured prediction to no-regret online
learning,'' in \emph{Proc. AISTATS}, 2011.

\bibitem{kirkpatrick2017}
J.~Kirkpatrick \emph{et al.},
``Overcoming catastrophic forgetting in neural networks,'' \emph{PNAS},
vol.~114, no.~13, pp.~3521--3526, 2017.

\bibitem{lifhoiem2017}
Z.~Li and D.~Hoiem,
``Learning without forgetting,'' \emph{IEEE TPAMI}, vol.~40, no.~12,
pp.~2935--2947, 2018.

\bibitem{mecklenburg2024}
L.~Mecklenburg \emph{et al.},
``Injecting new knowledge into large language models via supervised fine-tuning,''
arXiv:2404.00213, 2024.

\bibitem{chu2025}
T.~Chu, Y.~Zhai, J.~Yang, S.~Tong, S.~Xie, D.~Schuurmans, Q.~V.~Le, S.~Levine, and
Y.~Ma,
``SFT memorizes, RL generalizes: a comparative study of foundation model
post-training,'' in \emph{Proc. ICML}, 2025, pp.~10818--10838.

\bibitem{airef}
R.~Agarwal, N.~Vieillard, Y.~Zhou, P.~Sta{\'n}czyk, S.~Ramos, M.~Geist, and O.~Bachem,
``On-policy distillation of language models: learning from self-generated mistakes,''
in \emph{Proc. ICLR}, 2024, arXiv:2306.13649.

\bibitem{rlsrazor2025}
I.~Shenfeld, J.~Pari, and P.~Agrawal,
``RL's razor: why online reinforcement learning forgets less,''
in \emph{Proc. ICLR}, 2026, arXiv:2509.04259.

\bibitem{tml2025}
K.~Lu and Thinking Machines Lab,
``On-policy distillation,'' Thinking Machines Lab: Connectionism,
DOI 10.64434/tml.20251026, 2025.

\bibitem{bai2022}
Y.~Bai \emph{et al.},
``Constitutional AI: harmlessness from AI feedback,'' arXiv:2212.08073, 2022.

\bibitem{snell2022}
C.~Snell, D.~Klein, and R.~Zhong,
``Learning by distilling context,'' arXiv:2209.15189, 2022.

\bibitem{lora2022}
E.~J.~Hu \emph{et al.},
``LoRA: low-rank adaptation of large language models,'' in \emph{Proc. ICLR}, 2022.

\bibitem{qlora2023}
T.~Dettmers, A.~Pagnoni, A.~Holtzman, and L.~Zettlemoyer,
``QLoRA: efficient finetuning of quantized LLMs,'' in \emph{Proc. NeurIPS}, 2023.

\bibitem{unsloth}
Unsloth AI, ``Unsloth: 2x faster free finetuning,'' https://unsloth.ai, 2026.

\end{thebibliography}

\end{document}